\documentclass[xcolor=table, aspectratio=1610]{beamer}

\usetheme[]{Montpellier}

\setbeamertemplate{itemize items}[square]
\setbeamertemplate{section in toc}[sections numbered]

\setbeamertemplate{frametitle}{
    \begin{centering}
        \small\textsc{\insertframetitle}\par
    \end{centering}
}
\definecolor{mygreen}{RGB}{118, 188, 33}
\definecolor{mylightgreen}{RGB}{144,238,144}

\setbeamercolor*{palette primary}{bg=mygreen!70, fg=white}
\setbeamercolor*{palette secondary}{bg=mygreen!1, fg=black}
\setbeamercolor*{palette tertiary}{bg=mygreen!1, fg=black}
\setbeamercolor*{palette quaternary}{bg=mygreen!85, fg=white}

\setbeamercolor{block title}{bg=mygreen, fg=white}
\setbeamercolor{block body}{bg=mylightgreen, fg=black}
\setbeamertemplate{title page}[default][rounded=false]
\usepackage{amsfonts}
\usepackage{amsmath}
\usepackage{amsthm}
\usepackage{amssymb}
\usepackage{mathrsfs}
\usepackage{tikz}
\usepackage{graphicx}
\usepackage{subcaption}
\usepackage{booktabs}
\usepackage{makecell}
\usepackage{multirow}
\usepackage{multicol}
\usepackage{fouriernc}
\hypersetup{
    colorlinks=true,
    linkcolor=black,
    urlcolor=blue,
    citecolor=blue
}
\usetikzlibrary{arrows,calc}
\usepackage{relsize}

\makeatletter
\usepackage{xcolor}
\definecolor{customtitlebg}{RGB}{230, 230, 250}
\definecolor{customauthorbg}{RGB}{245, 245, 220}
\definecolor{customdatebg}{RGB}{240, 255, 240}

\setbeamercolor{title}{bg=mygreen, fg=white}

\makeatletter
\defbeamertemplate*{title page}{custom}{
  \vbox{}
  \vfill
  \begin{centering}
    \begin{beamercolorbox}[sep=14pt,rounded=false, shadow=true, wd=\textwidth, center]{title}
      \usebeamerfont{title}\inserttitle\par%
      \ifx\insertsubtitle\@empty%
      \else%
        \vskip0.25em%
        {\usebeamerfont{subtitle}\usebeamercolor[fg]{subtitle}\insertsubtitle}%
        \par%
      \fi%
    \end{beamercolorbox}%
    \vskip1em\par%
    \begin{beamercolorbox}[sep=8pt,center]{author}
      \usebeamerfont{author}\insertauthor\\
      {
\vspace{80pt}
\vspace{\fill}
\footnotesize{Last Updated: \today}
}
    \end{beamercolorbox}
    \vskip1em\par%
    \vfill
  \end{centering}
}
\makeatother
\useinnertheme{rounded}

\setbeamertemplate{title page}[custom]

\title[Sierra Leone--Summary]{\textsf{Summary Data on Acid Cropland and Lime Requirements in Sierra Leone}} 
\author[GAIA]{GAIA\\ \href{www.acidsoils.africa}{www.acidsoils.africa} } 
\institute[CIMMYT] 

\begin{document}

\begin{frame}[plain]
\textsc{\titlepage}
\end{frame}

\begin{frame}
\frametitle{Outlook}
\tableofcontents 
\end{frame}

\section{Introduction}
\section{Acidic Cropland}
\subsection*{Top 10 Provinces in Sierra Leone with the highest acidic cropland}

\begin{frame}{Top 10 Provinces in Sierra Leone with the highest acidic cropland}
    \begin{figure}
    \includegraphics[width=0.95\textwidth]{/Users/bisrathaile/Library/CloudStorage/Dropbox/CIMMYT-WORK/Analysis/bisrat-analysis-workflows/gaia_exante_bisrat/gaia-ex-ante-wow-shared/data-output//figures_and_tables/SLE//SLE_plt_acidic_cropland.png}
    \end{figure}
\end{frame}

\subsection*{Top 10 Provinces in Sierra Leone with the highest acidic cropland (Percentage of total cropland)}
\begin{frame}{Top 10 Provinces in Sierra Leone with the highest acidic cropland (Percentage of total cropland)}
    \begin{figure}
    \includegraphics[width=0.95\textwidth]{/Users/bisrathaile/Library/CloudStorage/Dropbox/CIMMYT-WORK/Analysis/bisrat-analysis-workflows/gaia_exante_bisrat/gaia-ex-ante-wow-shared/data-output//figures_and_tables/SLE//SLE_plt_acidic_cropland_perc.png}
    \end{figure}
\end{frame}

\section*{Total crop area distribution (ha)}
\subsection*{Total crop area in the top-10 most acidic Provinces, not segregated by acidic cropland}
\begin{frame}{Total crop area in the top-10 most acidic Provinces, not segregated by acidic cropland}
    \begin{center}
        \includegraphics[width=0.95\textwidth]{/Users/bisrathaile/Library/CloudStorage/Dropbox/CIMMYT-WORK/Analysis/bisrat-analysis-workflows/gaia_exante_bisrat/gaia-ex-ante-wow-shared/data-output//figures_and_tables/SLE//SLE_plt_crop_area_acidic.png}
    \end{center}
\end{frame}

\begin{frame}
    \scriptsize
    \begin{table}[!h]
\centering\begingroup\fontsize{7}{9}\selectfont

\begin{tabular}{>{\raggedright\arraybackslash}p{1.8cm}rrrrrrrrr}
\toprule
Provinces & Cacao & Cassava & Groundnut & Maize & \makecell{Other-\\Cereal} & \makecell{Other-\\Legume} & \makecell{Other-\\Non-food} & \makecell{Other-\\RTBs} & Sweet potato\\
\midrule
\cellcolor{gray!10}{Eastern} & \cellcolor{gray!10}{31.11} & \cellcolor{gray!10}{70.67} & \cellcolor{gray!10}{15.67} & \cellcolor{gray!10}{7.12} & \cellcolor{gray!10}{7.40} & \cellcolor{gray!10}{0.09} & \cellcolor{gray!10}{8.49} & \cellcolor{gray!10}{-} & \cellcolor{gray!10}{8.83}\\
Northern & 2.58 & 159.57 & 61.30 & 15.89 & 35.48 & - & 2.07 & 0.01 & 19.47\\
\cellcolor{gray!10}{Southern} & \cellcolor{gray!10}{8.19} & \cellcolor{gray!10}{149.04} & \cellcolor{gray!10}{27.41} & \cellcolor{gray!10}{10.60} & \cellcolor{gray!10}{16.94} & \cellcolor{gray!10}{-} & \cellcolor{gray!10}{4.01} & \cellcolor{gray!10}{-} & \cellcolor{gray!10}{10.80}\\
Western & 0.21 & 5.40 & 1.54 & 0.35 & 0.64 & - & 0.08 & - & 0.86\\
\bottomrule
\end{tabular}
\endgroup{}
\end{table}

\end{frame}

\section{Yield Loss}
\subsection*{Yield loss for Maize and Cassava(\%) in Sierra Leone}
\begin{frame}{Yield loss due to soil acidity for Maize and Cassava in Sierra Leone}
    \begin{figure}
    \includegraphics[width=0.95\textwidth]{/Users/bisrathaile/Library/CloudStorage/Dropbox/CIMMYT-WORK/Analysis/bisrat-analysis-workflows/gaia_exante_bisrat/gaia-ex-ante-wow-shared/data-output//figures_and_tables/SLE//SLE_yield_loss.png}
    \end{figure}
\end{frame}

\section{Lime Requirement}
\subsection*{Lime requirements (t) for acid soil remediation}
\begin{frame} {Lime requirements (t) for acid soil remediation in Sierra Leone}
    \begin{figure}
        \centering
        \includegraphics[width=0.95\textwidth]{/Users/bisrathaile/Library/CloudStorage/Dropbox/CIMMYT-WORK/Analysis/bisrat-analysis-workflows/gaia_exante_bisrat/gaia-ex-ante-wow-shared/data-output//figures_and_tables/SLE//SLE_plt_lime_requirements.png}
    \end{figure}
\end{frame}

\subsection*{Lime requirements (t) for acid soil remediation (table)}
\begin{frame} {Lime requirements (t) for acid soil remediation in Sierra Leone (table)}
    \scriptsize
    \begin{table}[!h]
\centering\begingroup\fontsize{7}{9}\selectfont

\begin{tabular}{>{\raggedright\arraybackslash}p{1.8cm}rrrrrrrrr}
\toprule
Provinces & Cacao & Cassava & Groundnut & Maize & \makecell{Other-\\Cereal} & \makecell{Other-\\Legume} & \makecell{Other-\\Non-food} & Sweet potato & \makecell{Other-\\RTBs}\\
\midrule
\cellcolor{gray!10}{Eastern} & \cellcolor{gray!10}{0.01} & \cellcolor{gray!10}{0.38} & \cellcolor{gray!10}{22.39} & \cellcolor{gray!10}{8.05} & \cellcolor{gray!10}{8.28} & \cellcolor{gray!10}{0.1} & \cellcolor{gray!10}{2.50} & \cellcolor{gray!10}{10.00} & \cellcolor{gray!10}{-}\\
Northern & 0.09 & 8.71 & 87.57 & 16.99 & 37.47 & - & 0.44 & 19.98 & 0.02\\
\cellcolor{gray!10}{Southern} & \cellcolor{gray!10}{0.03} & \cellcolor{gray!10}{1.33} & \cellcolor{gray!10}{39.16} & \cellcolor{gray!10}{10.38} & \cellcolor{gray!10}{16.55} & \cellcolor{gray!10}{0.0} & \cellcolor{gray!10}{0.44} & \cellcolor{gray!10}{10.75} & \cellcolor{gray!10}{-}\\
Western & - & - & 2.20 & 0.22 & 0.35 & - & - & 0.67 & -\\
\bottomrule
\end{tabular}
\endgroup{}
\end{table}

\end{frame}

\subsection*{Lime rate for Maize and Cassava in Sierra Leone}
\begin{frame} {Lime rate for Maize and Cassava in Sierra Leone}
    \begin{figure}
        \centering
        \includegraphics[width=0.98\textwidth]{/Users/bisrathaile/Library/CloudStorage/Dropbox/CIMMYT-WORK/Analysis/bisrat-analysis-workflows/gaia_exante_bisrat/gaia-ex-ante-wow-shared/data-output//figures_and_tables/SLE//SLE_lime_rate.png}
    \end{figure}
\end{frame}

\section{Yield Response}
\subsection*{Yield response to lime application in Sierra Leone}
\begin{frame} {Yield response to lime application in Sierra Leone}
    \begin{figure}
        \centering
        \includegraphics[width=0.98\textwidth]{/Users/bisrathaile/Library/CloudStorage/Dropbox/CIMMYT-WORK/Analysis/bisrat-analysis-workflows/gaia_exante_bisrat/gaia-ex-ante-wow-shared/data-output//figures_and_tables/SLE//SLE_yield_response.png}
    \end{figure}
\end{frame}

\section{Additional Production}
\subsection*{Potential for additional production (tone) in top 10 acidic Provinces in Sierra Leone}
\begin{frame}{Potential for additional production (tone) in top 10 acidic Provinces in Sierra Leone}
    \begin{figure}
    \includegraphics[width=0.98\textwidth, height=0.85\textheight]{/Users/bisrathaile/Library/CloudStorage/Dropbox/CIMMYT-WORK/Analysis/bisrat-analysis-workflows/gaia_exante_bisrat/gaia-ex-ante-wow-shared/data-output//figures_and_tables/SLE//SLE_plt_additional_production.png}
    \end{figure}
\end{frame}

\subsection*{Potential for additional production value (US\$) in top 10 acidic Provinces in Sierra Leone}
\begin{frame}{Potential for additional production value (US\$) in top 10 acidic Provinces in Sierra Leone}
    \begin{figure}
    \includegraphics[width=0.98\textwidth, height=0.85\textheight]{/Users/bisrathaile/Library/CloudStorage/Dropbox/CIMMYT-WORK/Analysis/bisrat-analysis-workflows/gaia_exante_bisrat/gaia-ex-ante-wow-shared/data-output//figures_and_tables/SLE//SLE_plt_additional_production_value.png}
    \end{figure}
\end{frame}

\section{ROI for Lime Application}
\subsection*{Return on investment for lime application in Sierra Leone}
\begin{frame} {Return on investment for lime application in Sierra Leone}
    \begin{figure}
        \centering
        \includegraphics[width=0.98\textwidth]{/Users/bisrathaile/Library/CloudStorage/Dropbox/CIMMYT-WORK/Analysis/bisrat-analysis-workflows/gaia_exante_bisrat/gaia-ex-ante-wow-shared/data-output//figures_and_tables/SLE//SLE_roi.png}
    \end{figure}
\end{frame}

\section{Profitability of Lime Application}
\subsection*{Profitability of lime application in Sierra Leone}
\begin{frame} {Profitability of lime application in Sierra Leone}
    \begin{figure}
        \centering
        \includegraphics[width=0.98\textwidth]{/Users/bisrathaile/Library/CloudStorage/Dropbox/CIMMYT-WORK/Analysis/bisrat-analysis-workflows/gaia_exante_bisrat/gaia-ex-ante-wow-shared/data-output//figures_and_tables/SLE//SLE_profit.png}
    \end{figure}
\end{frame}
\begin{frame}
\frametitle{Main assumptions and parameters}
\small
\begin{itemize}

    \item Acidic cropland is defined as having acidity saturation levels above 10\% of the Effective Cation Exchange Capacity (ECEC).
    \item Crop yields are unaffected by soil acidity until a specific acidity saturation threshold Target Acidity Saturation(TAS) is reached; beyond this threshold, yields decrease linearly, eventually reaching zero at high acidity levels.
    \item The TAS that separates maximum yield from yield decline is used as a critical value for making liming recommendations.
    \item Actual yields in acid tropical soils are a fraction of the potential yields in non-acidic soils, determined by the estimated yield loss due to acidity.
    \item Economic benefits from liming are calculated by multiplying the additional yield after soil remediation by median crop prices from sub-Saharan Africa (2016-2020 FAOSTAT data). The cost of liming is calculated using estimated lime application rates and a fixed lime price of 100 US\$/ton.
    \item The benefits of liming are assessed only for the year of application due to a lack of robust, large-scale data on long-term effects in sub-Saharan Africa.
\end{itemize}
For details, visit \href{www.acidsoils.africa}{www.acidsoils.africa}.
\end{frame}

\end{document}
